\section{Model Komunikasi Aman}

Model dasar komunikasi aman melibatkan beberapa pihak utama:

\begin{itemize}
\item \textbf{Pengirim (Sender/Alice):} Pihak yang mengirim pesan.
\item \textbf{Penerima (Receiver/Bob):} Pihak yang menerima pesan.
\item \textbf{Penyerang (Attacker/Eve):} Pihak yang mencoba menyadap atau memodifikasi pesan tanpa izin.
\item \textbf{Sumber tepercaya (Trusted Authority):} Pihak yang mendistribusikan kunci atau sertifikat.
\end{itemize}

\begin{figure}[!htbp]
\centering
\begin{tikzpicture}[
  node distance=1.2cm,
  box/.style={rectangle, draw, minimum width=1.8cm, minimum height=0.8cm, align=center}
]
\node[box, fill=blue!15] (alice) {Alice\\Pengirim};
\node[box, fill=green!15, right=3.5cm of alice] (bob) {Bob\\Penerima};
\node[box, fill=red!15, below=1.2cm of $(alice)!0.5!(bob)$] (eve) {Eve\\Penyerang};

\draw[-Stealth, thick] (alice) -- node[above, font=\scriptsize] {ciphertext} (bob);
\draw[-Stealth, red, dashed] (eve.north) -- ($(alice)!0.5!(bob)$);
\end{tikzpicture}
\caption{Model komunikasi kriptografi: Alice mengirim pesan terenkripsi ke Bob; Eve menyadap saluran.}
\label{fig:model-komunikasi}
\end{figure}

Dalam model klasik Shannon \cite{shannon1949}, Alice ingin mengirim pesan rahasia ke Bob melalui saluran yang tidak aman. Alice mengenkripsi pesan (plaintext) menjadi ciphertext menggunakan kunci, kemudian Bob mendekripsi ciphertext menjadi plaintext menggunakan kunci yang sama atau terkait. Eve yang menyadap saluran hanya melihat ciphertext dan tidak dapat memperoleh informasi bermakna tanpa kunci \cite{stinson2018}.
